% ===================================================================
%  CADRO N.º 9 — A montaña e a vila balnearia. — O bosque e a caza.
%  Traduction 2026 d'apres texto/io, controlee sur texto/fr, texto/es
%  et texto/pt. Consigne : voir texto/gl/10-cadro-01.tex.
%
%  LE TABLEAU DE LA MONTAGNE DONNE LE MEME RESULTAT QUE DANS LES
%  QUATRE COLONNES PRECEDENTES : le relief garde ses mots. « Penedo »
%  le rocher, « fraga » la futaie, « ribeira » la rive, « regato » le
%  ruisseau, « outeiro » la colline, « fervenza » la cascade,
%  « brañas » les paturages d'altitude, « toxo » l'ajonc. C'est le
%  cinquieme temoignage, apres l'ukrainien, le basque, le roumain et
%  l'irlandais.
%
%  MAIS LA CURE THERMALE, ELLE, EST INTERNATIONALE DES TROIS COTES :
%  « estación termal », « sanatorio », « augas minerais »,
%  « funicular », « viaduto », « casino », « villa ». Le meme partage
%  qu'a la maladie du tableau 4, et dans la meme page que la montagne
%  qui ne l'a pas — ce qui montre une derniere fois que le critere ne
%  tient pas au sujet du tableau mais au lieu ou le mot s'apprend.
% ===================================================================

%%K t09-apar-1 apar
\VUsaut{4.0mm}
\VUtitre{15.0pt}{60}{CADRO N.\textsuperscript{o} 9}
\VUsaut{3.0mm}

%%K t09-tit-1 sub
\VUcentre{12.0pt}{0}{\textit{Primeira escena.}}
\VUsaut{3.0mm}

%%K t09-tit-2 sub
\VUpk{11.4pt}{40}{A montaña. --- A vila balnearia.}
\VUsaut{3.0mm}

%%K t09-c1-01-1 p
1. --- Levo oito días en Pratz. Non podo expresar toda a admiración que sentín cando
estiven diante da marabillosa \VUgras{serra} \textsuperscript{(1)} dos Pireneos, cuxa
\VUgras{crista} \textsuperscript{(2)} se recorta no ceo azul coma un festón xigante.

%%K t09-c1-02-1 p
2. --- Non imaxinaba que os \VUgras{cumes} dos Pireneos \textsuperscript{(3)}
acadasen unha \VUgras{altura} tan grande, e quedei abraiado diante dos fermosísimos
\VUgras{glaciares} \textsuperscript{(5)} e dos \VUgras{picos} nevados
\textsuperscript{(4)}, polos que rolan \VUgras{avalanchas} tan terroríficas.

%%K t09-tit-3 sub
\VUsaut{3.0mm}
\VUpk{10.2pt}{40}{Paisaxe de montaña.}
\VUsaut{3.0mm}

%%K t09-c1-03-1 p
3. --- Xa ao día seguinte da miña chegada fun ao vello \VUgras{volcán}
\textsuperscript{(6)} apagado que está preto de Pratz. Nas bordas áridas e espidas do
\VUgras{cráter} vense aínda ben os vestixios do paso da \VUgras{lava}, que correu,
hai varios séculos, pola \VUgras{ladeira da montaña} \textsuperscript{(7)}. Conseguín
atopar, nunha das \VUgras{costas}, algúns anacos desa lava.

%%K t09-c1-04-1 p
4. --- Pratz está situada na fronteira, e a \VUgras{fortaleza}
\textsuperscript{(8)} que defende o \VUgras{desfiladeiro} \textsuperscript{(27)} que
separa Francia de España lembra por completo eses vellos castelos da ribeira do Rin
que semellan axexar unha \VUgras{presa} desde o alto do seu penedo.

%%K t09-c1-05-1 p
5. --- Unha enorme \VUgras{aguia} \textsuperscript{(9)}, que ten o seu niño non lonxe
do \VUgras{forte} \textsuperscript{(8)}, atrae moitas veces a miña atención. Esa
\VUgras{ave de rapina}, cuxo \VUgras{bico} é forte, leva a miúdo aos seus fillos un
año ou un cabrito que sostén coas súas \VUgras{unllas.} É tanta a grandeza das súas
\VUgras{ás} que pode planear moito tempo coa súa pesada carga.

%%K t09-tit-4 sub
\VUsaut{3.0mm}
\VUpk{10.2pt}{40}{O excursionista.}
\VUsaut{3.0mm}

%%K t09-c1-06-1 p
6. --- Alí o paseo máis agradable é a excursión á cascada \textsuperscript{(10)} de
«Cau». A \VUgras{fervenza} cae en remuíños e despois desaparece como un belo
\VUgras{regato} nun \VUgras{bosque de abetos} \textsuperscript{(11)} situado ao oeste
da vila. Subimos por ese bosque ata o \VUgras{observatorio meteorolóxico}
\textsuperscript{(12)} e, desde o alto do \VUgras{miradoiro}, albiscamos todo o
\VUgras{panorama} da rexión. A \VUgras{paisaxe} é marabillosa, e a
\VUgras{natureza} semella prodigar á terra todos os tesouros de que dispón.

%%K t09-c1-07-1 p
7. --- Para baixar, collemos o \VUgras{funicular} \textsuperscript{(13)} e paramos
nunha pequena \VUgras{estación} preto das \VUgras{ruínas} \textsuperscript{(14)} dun
vello \VUgras{castelo feudal} habitado noutrora polos señores de Pratz.

%%K t09-c1-08-1 p
8. --- Pobre castelo! Que pensará ao ver abaixo a coqueta \VUgras{vila balnearia}
\textsuperscript{(15)} que hoxe é unha \VUgras{estación termal} de moda?

%%K t09-c1-08-2 p
O \VUgras{clima} é tan agradable, as \VUgras{augas minerais} tan eficaces, que a
\VUgras{cura} que se segue no \VUgras{sanatorio} \textsuperscript{(17)} é un
verdadeiro pracer.

%%K t09-tit-5 sub
\VUsaut{3.0mm}
\VUpk{10.2pt}{40}{Unha agradable vila termal.}
\VUsaut{3.0mm}

%%K t09-c1-09-1 p
9. --- Ademais, fíxose todo para facer agradable a estadía. Construíuse un gran
\VUgras{viaduto} \textsuperscript{(16)} para permitir o ferrocarril; o
\VUgras{parque} \textsuperscript{(18)} do \VUgras{establecemento termal}, no que se
dan \VUgras{concertos}, está disposto dun xeito encantador. Construíronse varias
graciosas «\VUgras{villas}» \textsuperscript{(19)} arredor do casino
\textsuperscript{(21)}, e unha \VUgras{estrada} \textsuperscript{(22)} moi ben
mantida facilita moito as comunicacións.

%%K t09-c1-10-1 p
10. --- Un simple \VUgras{noiro} \textsuperscript{(23)} separa o \VUgras{camiño}
\textsuperscript{(20)} do \VUgras{val} \textsuperscript{(25)} propiamente dito, no
fondo do cal xace un gracioso \VUgras{lago} pequeno \textsuperscript{(26)} que
alimenta un \VUgras{regato} \textsuperscript{(24)} límpido.

%%K t09-c1-11-1 p
11. --- Do outro lado do val explótase unha \VUgras{mina de ferro}
\textsuperscript{(28)}. Fun varias veces ver os \VUgras{mineiros} baixando ao
\VUgras{pozo}, e mesmo entrei na \VUgras{galería} na que pasan o día, extraendo o
\VUgras{mineral} que contén o \VUgras{metal.}

%%K t09-c1-12-1 p
12. --- Construíuse unha importante \VUgras{fundición} preto da mina para aproveitar
axiña as riquezas que dela se extraen. O \VUgras{forno alto} \textsuperscript{(29)},
quentado co \VUgras{carbón de pedra} que abonda aquí, permite obter ferro
\VUgras{fundido} axiña e barato.

%%K t09-c1-13-1 p
13. --- Non lonxe da mina escávase unha \VUgras{canteira} \textsuperscript{(30)}. Nin
\VUgras{granito,} nin \VUgras{pórfido,} nin sequera \VUgras{gres} corta o vello
\VUgras{canteiro} co seu \VUgras{pico,} senón simple \VUgras{pedra de cantaría}
para a construción das casas.

%%K t09-c1-14-1 p
14. --- Un dos máis belos adornos desa terra son os \VUgras{abetos}
\textsuperscript{(31)}. En todas as partes, no fondo dunha \VUgras{cavorca}
\textsuperscript{(32)}, na ribeira dun \VUgras{torrente} \textsuperscript{(33)},
dunha \VUgras{sima} \textsuperscript{(34)} ou dun precipicio, as súas agullas
delgadas botan matices verdes.

%%K t09-tit-6 sub
\VUsaut{3.0mm}
\VUpk{10.2pt}{40}{Andar a pé.}
\VUsaut{3.0mm}

%%K t09-c1-15-1 p
15. --- Aproveito as miñas vacacións para pasear \VUgras{moito}. Os \VUgras{camiños}
\textsuperscript{(35)} que serpean pola montaña permiten percorrer, sen reparar na
distancia, varios \VUgras{quilómetros} e a miúdo varias \VUgras{decenas de
quilómetros}.

%%K t09-c1-15-2 p
\VUgras{Marcos} \textsuperscript{(36)} e \VUgras{postes indicadores}
\textsuperscript{(37)} sinalan os camiños, nos que moitas veces se atopa un mozo
\VUgras{montañés} \textsuperscript{(38)} coa súa \VUgras{alforxa}
\textsuperscript{(40)} ao lado, levando unha \VUgras{marmota}
\textsuperscript{(39)}, ou algún \VUgras{cigano} \textsuperscript{(41)} cuxa
\VUgras{chaqueta de pel} \textsuperscript{(42)} contrasta estrañamente co vello
\VUgras{capote} \textsuperscript{(43)} esfarrapado. Leva cunha \VUgras{cadea un oso}
\textsuperscript{(44)} amestrado.

%%K t09-tit-7 sub
\VUpk{10.2pt}{40}{Ascensión ás montañas.}
\VUsaut{3.0mm}

%%K t09-c1-16-1 p
16. --- Case cada día vou cun \VUgras{guía} \textsuperscript{(48)} practicar a
\VUgras{ascensión}, e estou moi orgulloso cando, coa \VUgras{mochila}
\textsuperscript{(47)} ás costas e o «\VUgras{alpenstock}» na man, coma un verdadeiro
\VUgras{turista} \textsuperscript{(46)}, asalto os \VUgras{penedos}
\textsuperscript{(45)}.

%%K t09-c1-17-1 p
17. --- Desde o alto dunha montaña, os habitantes da chaira semellan non máis grosos
ca formigas; un \VUgras{rabaño de ovellas} \textsuperscript{(50)} que pace nun
\VUgras{pasteiro} \textsuperscript{(49)} semella un xoguete para nenos. Un
\VUgras{piñeiro} \textsuperscript{(51)} ou mesmo unha \VUgras{casiña de madeira}
\textsuperscript{(52)} apenas parece unha mancha sobre a verdura.

%%K t09-c1-18-1 p
18. --- Durante esas excursións atopamos a miúdo no \VUgras{carreiro}
\textsuperscript{(53)} un \VUgras{cazador de rebezos} \textsuperscript{(54)} que leva
ao ombreiro a besta que acaba de matar.

%%K t09-c1-19-1 p
19. --- Puxen no meu herbario unha póla moi notable de \VUgras{fento}
\textsuperscript{(55)} e unha bela pequena póla rosada de \VUgras{uz}
\textsuperscript{(57)} que colleitei na entrada dunha pequena \VUgras{gruta}
\textsuperscript{(56)} no meu último paseo.

%%K t09-tit-8 sub
\VUsaut{3.0mm}
\VUcentre{12.0pt}{0}{\textit{Segunda escena.}}
\VUsaut{3.0mm}

%%K t09-tit-9 sub
\VUpk{11.4pt}{40}{O bosque. --- A caza.}
\VUsaut{3.0mm}

%%K t09-c2-01-1 p
1. --- Onte pasei un día delicioso no bosque. A dilixencia que reina entre os animais
e os \VUgras{insectos} \textsuperscript{(5)} que o habitan interesoume moito.

%%K t09-c2-01-2 p
A \VUgras{araña} \textsuperscript{(1)} que tece a súa \VUgras{tea}
\textsuperscript{(2)} entre dúas pólas para apañar a \VUgras{mosca}
\textsuperscript{(3)} que pasa, o \VUgras{caracol} \textsuperscript{(4)} que leva a
custo a súa cuncha, o vacaloura que anda á procura da súa presa, a formiga dilixente
moi atarefada preto do \VUgras{formigueiro} \textsuperscript{(6)}, todos eses
pequenos iban, viñan e traballaban cunha dilixencia extraordinaria.

%%K t09-c2-02-1 p
2. --- Unha \VUgras{serpe} \textsuperscript{(7)} que a primeira vista tomei por unha
\VUgras{víbora}, pero que non era outra cousa ca unha simple \VUgras{cóbrega}, estaba
acochada baixo o tronco dunha árbore e de cando en vez sacaba a súa cabeciña delgada,
que sempre semellaba disposta a morder. Diante de min unha grosa \VUgras{lesma}
\textsuperscript{(8)} reptaba pesadamente despois de devorar unha das saborosas
\VUgras{amorodiñas} \textsuperscript{(11)} que medran alí en abundancia.

%%K t09-c2-03-1 p
3. --- O chan estaba tapizado de \VUgras{flores do bosque}: \VUgras{primaveras}
\textsuperscript{(9)}, \VUgras{pervincas} \textsuperscript{(10)}, e moitas outras
plantas que eu non coñecía florecían entre as \VUgras{moutas de herba}
\textsuperscript{(12)}.

%%K t09-c2-04-1 p
4. --- Nesa rexión a \VUgras{trufa} é case tan común coma o \VUgras{cogomelo}
\textsuperscript{(14)}, e un \VUgras{porquiño} \textsuperscript{(13)} que pertencía a
un gardabosques buscaba con gula se conseguiría atopar algunhas delas.

%%K t09-tit-10 sub
\VUsaut{3.0mm}
\VUpk{10.2pt}{40}{Caza furtiva.}
\VUsaut{3.0mm}

%%K t09-c2-05-1 p
5. --- Asistín ao \VUgras{final} dunha escena que me divertiu moito. Coa axuda do
olfacto do seu \VUgras{can basset} \textsuperscript{(15)}, o \VUgras{gardabosques}
\textsuperscript{(16)} acaba de sorprender un \VUgras{cazador furtivo}
\textsuperscript{(22)} no intre en que este se preparaba para levar a \VUgras{caza}
\textsuperscript{(17)} que matara. Preferindo abandonar a súa presa a deixarse
prender, o cazador furtivo fuxira, e a \VUgras{lebre} \textsuperscript{(18)}, a
\VUgras{cerva} \textsuperscript{(19)}, o \VUgras{corzo} \textsuperscript{(20)} e o
\VUgras{xabaril} \textsuperscript{(21)} quedaron na herba, alegrando o gardabosques.

%%K t09-c2-06-1 p
6. --- A variedade das árbores que ten o bosque é abraiante. As \VUgras{faias}
\textsuperscript{(23)} e os \VUgras{bidueiros} \textsuperscript{(26)}, os
\VUgras{piñeiros}, que dan a \VUgras{resina} \textsuperscript{(27)}, os
\VUgras{carballos} \textsuperscript{(29)} e moitos outros mesturan a súa ramallada.

%%K t09-c2-06-2 p
A \VUgras{hedra} \textsuperscript{(24)}, que se agarra ao tronco das árbores altas,
colora a \VUgras{fraga} \textsuperscript{(25)} cun verdor brillante que se une
agradablemente ao matiz máis claro e verde pálido do \VUgras{musgo}
\textsuperscript{(31)}, no que o \VUgras{peto verde} \textsuperscript{(30)} busca
insectos.

%%K t09-tit-11 sub
\VUsaut{3.0mm}
\VUpk{10.2pt}{40}{Paisaxe do bosque.}
\VUsaut{3.0mm}

%%K t09-c2-07-1 p
7. --- Os vellos carballos encantáronme. As súas grosas \VUgras{raíces}
\textsuperscript{(32)} retortas, medio saídas do chan, danlles un espléndido aspecto de
forza e vigor, e pregúntome como o \VUgras{propietario} \textsuperscript{(35)} pode
resignarse a facelos abater e entregalos á \VUgras{serra} \textsuperscript{(34)} do
\VUgras{leñador} \textsuperscript{(33)}. Os pobres \VUgras{troncos}
\textsuperscript{(36)}, espidos da súa \VUgras{cortiza} \textsuperscript{(37)} ou
fendidos pola \VUgras{machada} \textsuperscript{(38)} do \VUgras{xornaleiro}
\textsuperscript{(39)}, danme pena.

%%K t09-c2-08-1 p
8. --- Preto de alí, un vello \VUgras{carboeiro} \textsuperscript{(41)} preparou coa
súa \VUgras{pa} un \VUgras{montón} \textsuperscript{(42)} de carbón, e unha vella
levaba un \VUgras{feixe} \textsuperscript{(40)} de \VUgras{leña morta} feito coas
pólas das árbores abatidas.

%%K t09-tit-12 sub
\VUsaut{3.0mm}
\VUpk{10.2pt}{40}{A caza.}
\VUsaut{3.0mm}

%%K t09-c2-09-1 p
9. --- Quixen ir descansar uns intres na \VUgras{casa do gardabosques}
\textsuperscript{(46)}, pero, cando cheguei preto da \VUgras{lagoa}
\textsuperscript{(43)}, na que nada un belo \VUgras{cisne} \textsuperscript{(45)},
reparei en que unha \VUgras{inundación} recente \textsuperscript{(44)} mudara o
camiño nun verdadeiro \VUgras{lameiro}. Tiven que dar a volta e, ao pasar preto do
\VUgras{cedro} \textsuperscript{(49)}, dos \VUgras{chopos} \textsuperscript{(48)} e do
\VUgras{ulmeiro} \textsuperscript{(47)} que están á beira do bosque, vin saír dun
\VUgras{souto} \textsuperscript{(54)} un fermosísimo \VUgras{cervo}
\textsuperscript{(50)}, perseguido por unha teimosa \VUgras{matilla}
\textsuperscript{(51)} que tamén azuzaba o \VUgras{corno} \textsuperscript{(53)} dos
\VUgras{cazadores a cabalo} \textsuperscript{(52)}. A pobre besta estaba de todo
esgotada. Veu caer morrendo a uns pasos de min.

%%K t09-c2-10-1 p
10. --- O fillo do gardabosques, a quen atraera o ruído da \VUgras{montaría}, saíu da
casa e volvemos os dous ao bosque. Vira unha \VUgras{pega} \textsuperscript{(56)}
nunha póla dun vello carballo. Coidaba que o seu niño debía de estar agochado na
\VUgras{follaxe} \textsuperscript{(58)} e quixo apañar o niño.

%%K t09-c2-10-2 p
Áxil coma un \VUgras{esquío} \textsuperscript{(61)}, subiu de \VUgras{póla}
\textsuperscript{(55)} en póla, pero non atopou nada e só conseguiu facer botar a voar
unha vella \VUgras{corvella} \textsuperscript{(60)} pousada nun \VUgras{ramallo}
\textsuperscript{(57)}.
\VUsaut{4.0mm}
\VUfilet{22mm}
